% =====================================================================
% overleaf/01_introduction.tex — LaTeX rendering of en/01_introduction.md v0.17 (11 September 2026)
% Dernière mise à jour : 2026-09-11
% Chapter 1 of the paper. The related work is a \section of its own: overleaf/02_related_work.tex,
% to be inserted right after this file. Cross-references use \label / \ref, so no number is hard-coded.
% Ready to paste into the Overleaf project (acmart + natbib: \citep, \citet).
% Derived file: regenerate after every validated pass of en/01_introduction.md.
% Citation keys are those of docs/paper/sources/references.bib (checklist: article/CITATIONS.md).
% Preamble additions:
%   \newcommand{\todo}[1]{\textbf{[#1]}}   % placeholder figures, filled from methode/experience_plan/experiments.yaml
% Vocabulary: Level 0--3 = our ablation levels; exploratory / robust / transferable = SILICA certification steps.
% Comments "% source:" give the repository file a plain-text figure is recomputed from.
% =====================================================================
% =====================================================================

\section{Introduction}
\label{sec:intro}

\subsection{Context: multimodal mobility, user-centric transport and the fidelity requirement of planning}

Understanding individual mobility behaviour in multimodal transport systems is essential to advance travel planning and to design efficient networks. The challenge is not only to reproduce aggregate flows: it is to build a simulation framework able to support the study of individual behaviour, the testing of fine-grained behavioural hypotheses and the design of more sustainable, personalised mobility policies. For this, the simulated population must also reproduce, at least in aggregate, the behaviour recorded by certified household travel surveys \citep{tisseo2023emc2}.

The rise of personalised multimodal transport systems reflects a transition toward user-centric mobility, in which public transit, walking, cycling and shared services are combined to offer journeys tailored to individual preferences \citep{smith1995transims,horni2016matsim}. Representing such systems in simulation is particularly complex, because mobility choices are shaped by heterogeneous socioeconomic profiles, accessibility constraints and lived experiences \citep{grignard2018impact}. As \citet{oberoi2024personalisation} underlines, personalisation within Mobility-as-a-Service requires dealing with the heterogeneity of users, their preferences, their life contexts and their accessibility needs, which remains a major obstacle to realistic transport models. Modelling must therefore combine flexibility and realism to capture heterogeneous, adaptive travel decisions at population scale.

Current transport simulations struggle with this complexity and suffer from a shortage of calibration data: agent-based models require extensive local data, which hinders their adoption \citep{liu2024toward}. Detailed multimodal trajectory and itinerary datasets, indispensable to grasp behavioural nuances, are extremely scarce \citep{fourez2025transport}. They exist only for a restricted number of large metropolitan areas, Tokyo for instance, and their coverage is very uneven from one city to another \citep{feng2024agentmove}. Moreover, integrating multiple external sources imposes heavy manual rule engineering, producing rigid models with fixed heuristics, costly to scale and helpless in unforeseen conditions.

Over the last three years, large language models (LLMs) have entered agent-based simulation as decision engines. Generative agents endowed with natural-language reasoning, episodic memory and planning were introduced by \citet{park2023generative}. They have since been proposed for travel demand modelling \citep{liu2024toward}, for large-scale urban simulation \citep{bougie2025citysim}, and for route replanning under road disruptions on the GAMA platform \citep{alves2026evaluating}. This work builds on the GAMA--OpenTripPlanner--LLM architecture of \citet{vu2025modeling}, applied here to the Toulouse metropolitan area.

The appeal of these agents answers the two difficulties above. They promise to represent decisions that no tabulated variable captures: a strike, a heatwave, a news item about safety in the metro. And they promise to do so without the calibration data that classical approaches require: a model that reasons from a persona and a set of itineraries seems to bypass both the scarcity of trajectories and the brittleness of hand-written rules.

Whether it does is an empirical question, and the published answers are partial. Most evaluations compare LLM agents with rule-based agents on scenario plausibility, or on aggregate patterns of their own making. One recent system, \citep{lammer2026gta}, does confront a simulated modal split with a national travel survey; but the only reference it gives itself is the observed modal split of other regions, and on Berlin one of them, Hamburg, is more than twice as close to the target as the simulation itself. What is missing is therefore not the measurement. It is a protocol under which a measurement becomes attributable: a floor beneath the agents, a calibrated model above them, and the same information given to every model compared. This paper supplies that protocol and applies it against a certified survey.

\subsection{The gap: linguistic plausibility is not distributional fidelity}
\label{sec:gap}

An LLM agent produces fluent, individually reasonable justifications for its mode choices. Fluency is not evidence. A population of such agents may be plausible one agent at a time and wrong in aggregate, because each agent draws on the same pre-training prior rather than on the behaviour of the territory it is meant to represent. \citet{meister2024benchmarking} observe the same dissociation at the level of a single model: it may know a distribution and still be unable to behave according to it.

SILICA \citep{bintareaf2026silica} gives this concern an operational form. It runs twelve open-weight models through five game-theoretic environments with published human anchors, 9,115 model runs plus 150 runs of classical baselines, and certifies each claim on a three-step ladder: \textbf{exploratory} if it holds at baseline in at least one model; \textbf{robust} if it also survives a library of perturbations that leave the rules unchanged, such as reordering the options or resetting memory, across at least three model families; \textbf{transferable} if, in addition, it matches the human anchor by an equivalence test and the transcripts are consistent with the human mechanism. No finding in SILICA reaches the transferable step; convention formation alone reaches the robust step. The author is explicit about the scope: small open-weight models on a consumer graphics card, so nothing is established about frontier models; anchors that are aggregate human values, four of the seven resting on a meta-analytic synthesis and the other three on a single experiment, obtained overwhelmingly on Western university samples; a single-author analysis.

Two things follow for urban mobility, and neither is a conclusion. First, the domain is different. Mode choice is an individual decision conditioned on a rich context, distance, schedules, vehicle availability, not a repeated interaction with strategic payoffs; and the reference is a certified survey of the same territory, not a meta-analytic mean. Whether the ceiling SILICA observes on economic games appears on mode choice is an open question. Second, the alignment literature of Section~\ref{sec:alignment} shows that supervision moves distributions. Our own use of the survey is deliberate and narrow. We read its published modal shares, compare them with the split our agents produce, and rewrite the prompt by hand to close the gap; the prompt carries behavioural guidance in words, never a figure from the survey, a target share, or an observed choice. The question is whether a city's own survey is enough to adapt an uncalibrated simulation to that city, or whether the gap is structural. Section~\ref{sec:contributions} phrases this as a single testable hypothesis.

We also make a distinction that organises the paper. A population of LLM agents can be asked two different things: to serve as a \textbf{statistical stand-in} for the population of a territory, reproducing its modal split; or to serve as an \textbf{adaptive system} that reacts to contexts no survey tabulates, a strike, a heatwave, a service breakdown, and carries the memory of a bad experience from one day to the next. The first use is judged by distributional fidelity. The second is judged by the robustness and direction of the reaction, because no survey follows the same individuals through such events. Conflating the two has produced both over-claims and unfair dismissals of generative agents. We evaluate them separately.

\subsection{Contributions and hypotheses}
\label{sec:contributions}

The study is conducted on the Toulouse metropolitan area, against the 2023 Cerema-certified household travel survey (EMC\textsuperscript{2} 2023) \citep{tisseo2023emc2}, which covers about 16,000 respondents over 453 municipalities. % source: rapport AUAT/CEREMA 2023 ; BIBLIOGRAPHIE.md §5

\paragraph{C1. An evaluation contract at informational parity.}
Our first contribution is methodological, and it is the one we consider most transferable to other evaluations of LLM agents. The contract has three rules. Each closes one way in which a comparison between a generative agent and a tabular model can be biased.

\begin{enumerate}
  \item \emph{Informational parity.} The LLM agent, the multinomial logit and the supervised oracle, a LightGBM model \citep{ke2017lightgbm} fitted on the survey microdata, receive the same 21 input variables, whose list is frozen in the repository. % source: scripts/progedo_logit/feature_spec.json
  Without this rule, any gap between models is confounded with a gap in information. Parity concerns inputs, not exposure: the oracle has been fitted on 31,279 survey trips % source: scripts/progedo_logit/mode_choice_policy_metrics.json, training.n_fit
  and the LLM has seen none. The oracle is therefore a reference ceiling, not a competitor.
  \item \emph{Like-for-like readings.} The agent does not pick one itinerary. Following the finding of \citet{meister2024benchmarking} that verbalised distributions align better than sampled ones, it is asked to spread 100\,\% over the itineraries presented, up to six per trip. % source: experience.yaml of the played experiments, max_candidats: 6
  The platform scores this answer as a probability mass per mode, compared with the renormalised probabilities of the tabular model. The agent's decision inside the simulation is a \textbf{draw} from that mass, seeded on the decision context, so that a population restores individual variability instead of a deterministic argmax; % source: mobility_llm/src/mobility_llm/mode_choice.py, draw_index
  the most probable mode is published as a secondary reading, against the tabular argmax. % source: scripts/synthesis/build.py, readings attendu / elu / brut
  A probabilistic model has two L1 errors, and only one of them is comparable to hard decisions.
  \item \emph{Renormalisation on the offer.} The tabular model predicts over four modes blindly; the agent chooses only among the itineraries the router actually returns. Every tabular prediction is restricted to the modes offered for that trip and renormalised to 100\,\%, an independence-of-irrelevant-alternatives assumption that we declare as a limit. Without this rule, the tabular model is credited for options that do not exist and the agent is penalised for modes it was never offered.
\end{enumerate}

The evaluation substrate is a sealed synthetic cohort of 1,000 personas, % source: data/population/population_1000_AAMAS/MANIFEST.yaml, seal 1 of 2026-09-02, sha256 f67b0777…
drawn by proportional allocation on residence ring $\times$ motorisation from an eqasim pool of 5,063 persons \citep{horl2021eqasim}, and controlled against the six margins for which the EMC\textsuperscript{2} survey provides a target, four published in the AUAT report and two recomputed from its microdata: age class, occupation, motorisation at person and household level, residence ring, and ring $\times$ motorisation. All six are judged conformant by the control procedure. % source: data/population/population_1000_AAMAS/CONTROLE.md
Their activity chains produce 2,693 trips on the evaluated day, of which 2,645 are evaluable. % source: compteurs.json of the executions of 2026-09-07
The tabular models are additionally characterised on the sealed survey test set of 13,045 trips, split by household. % source: mode_choice_policy_metrics.json, test.n_rows, split.by = hh_id
GTA aside, no LLM-mobility evaluation of Section~\ref{sec:agents} confronts an agent population with an observed modal split at all; and none, GTA included, does so with a floor beneath the agents, a calibrated reference above them and informational parity between the two.

\paragraph{C2. An ablation in four levels, and two hypotheses on the distributional ceiling.}
We compare, on the same sealed substrate, four levels of increasing information. Level~0: floors and physical heuristics, a uniform draw over the offer, the majority mode, the shortest itinerary. Level~1: the bare LLM, a neutral instruction and no behavioural guidance. Level~2: the same LLM with the hand-calibrated prompt of Section~\ref{sec:gap}, adjusted on a calibration population drawn separately from the sealed cohort. Level~3: the tabular references, multinomial logit and LightGBM oracle. Several LLMs are evaluated, all at temperature~0. One hypothesis is pre-registered.

\textbf{H0, distributional ceiling.} Neither the bare nor the calibrated LLM reaches the tabular reference within the equivalence margin on the four-mode split. Bare LLM L1 error: \todo{xx.y | exp\_01a} and \todo{xx.y | exp\_01b}. Calibrated prompt: \todo{xx.y | exp\_02a}. Logit: \todo{xx.y | exp\_03a}. Oracle: \todo{xx.y | exp\_03b}. The calibrated figure is measured once on the sealed cohort, on individuals that were not used to adjust the prompt. H0 is refuted if the calibrated condition enters the equivalence margin of the tabular reference.

\paragraph{C3. Two regimes no survey tabulates, evaluated for robustness rather than fidelity.}

\begin{itemize}
  \item \emph{Five-day hysteresis after a service breakdown.} On day~2 of a five-day horizon, metro line~A fails at 17:00 and imposes a 45-minute delay; the network is nominal again from day~3. % source: experiments.yaml, exp_04a perturbation
  Three conditions are compared: an agent with a short-term memory buffer in the sense of \citet{park2023generative}, decaying at rate $\lambda$; the same agent without memory; the tabular oracle, structurally memoryless. No survey follows the same travellers through such an event, so no value is a target. What is testable is the order of the conditions on day~3, \todo{xx.y | exp\_04a vs exp\_04b vs exp\_04c}, the monotonic recovery over days 3 to 5, and the displacement of the curve when $\lambda$ is divided by three, \todo{xx.y | exp\_04d}. The hypothesis falls if the memoryless agent shows the same inertia, if $\lambda$ has no effect, or if recovery is not monotonic.
  \item \emph{Five real, dated local news events.} Five real, dated articles of the Toulouse local press, chosen before any model call, provide five events. % source: RAPPORT_SCENARIOS_QUALITATIFS_TOULOUSE.md §2 and §3
  Each acts through a distinct channel, and none is representable as a tabular variable: gusts of the Autan wind above 80\,km/h closing the city's parks, a weather and physical-safety channel at metropolitan scale; a bed-bug scare in the metro and buses, a perceived-risk channel at network scale; an open-ended strike of waste collectors in the historic centre, a sensory-nuisance channel at district scale; the street parade of La Machine with a pedestrianised centre, a closure-and-attraction channel on the centre and its radial axes; the launch of electric bike-sharing on the hilly districts, a supply channel at metropolitan scale. Each event is run under five conditions on the same personas: the nominal day, with no article; the raw article; a mode-neutral paraphrase, the same fact rewritten without any mention of a mode; a neutral control text, a real local article of similar length with no plausible link to mode choice, which measures the effect of adding any text at all; and the tabular oracle with the event encoded into its variables where the event admits such an encoding. For each event and each mode, the expected direction of the shift was written down before any model call, 20 signed predictions in all; we report the share of observed modal shifts that go in the expected direction \todo{xx.y | exp\_05}, with option order randomised for every request, a control made necessary by the order sensitivity SILICA measures \citep{bintareaf2026silica}. The raw-minus-paraphrase contrast isolates instruction following from reasoning; the contrast with the neutral control text separates the effect of the article's content from the effect of adding a text.
\end{itemize}

\subsection{Paper organisation}

Section~2 reviews the related work, from random utility models to generative agents in mobility simulation and to the distributional alignment of LLM populations. Section~3 defines the macro and micro metrics, the three rules of the contract, and the demographic control of the sealed cohort. Section~4 evaluates the bare LLMs and their inter-seed variability. Section~5 presents the four-level ablation, the calibrated condition and the tabular references on the sealed survey test set, and tests H0. Section~6 presents the two non-tabulated regimes, hysteresis and news events, under the pre-registered directional predictions. Section~7 discusses the limits, including the independence-of-irrelevant-alternatives assumption of the renormalisation and the exposure asymmetry between oracle and agents, and the implications for hybrid designs. Section~8 concludes.
